% Source: https://artofproblemsolving.com/community/c3h1609462_utilizing_base_10_algebraically
% Stable identity: aops-topic-1609462-q01
Let $m=23\,000\,000$ and $n=24\,000\,000$. The conditions mean
\[
5\cdot10^{m-1}\le A<6\cdot10^{m-1},\qquad 9\cdot10^{n-1}\le B<10^{n}.
\]
Multiplying,
\[
4.5\cdot10^{m+n-1}\le AB<6\cdot10^{m+n-1}.
\]

\textbf{1.} Since $10^{m+n-1}\le AB<10^{m+n}$, the product has $m+n=\boxed{47\,000\,000}$ digits.

\textbf{2.} $AB$ lies in $[4.5,6)\cdot10^{m+n-1}$, so its first digit is $4$ or $5$, and both occur: $A=5\cdot10^{m-1}$, $B=9\cdot10^{n-1}$ give $AB=45\cdot10^{m+n-2}$ (first digit $4$), while $A=59\cdot10^{m-2}$, $B=99\cdot10^{n-2}$ give $AB=5841\cdot10^{m+n-4}$ (first digit $5$). So the first digit is $\boxed{4\text{ or }5}$; it is not determined.

\textbf{3.} No. If the first digit is $4$, then $AB\in[4.5,5)\cdot10^{m+n-1}$ and the second digit is one of $5,6,7,8,9$; the examples $45\ldots$ and, say, $A=55\cdot10^{m-2}$, $B=9\cdot10^{n-1}$ giving $495\cdot10^{m+n-3}$, show that it varies. If the first digit is $5$ the second digit can be anything from $0$ to $9$. So the second digit is $\boxed{\text{not determined}}$.
